\bookchapter{บทที่ 14: เจาะลึกตัวจัดการรหัสผ่าน}

\chapterintro{ในยุคดิจิทัล รหัสผ่านคือกุญแจบ้านของคุณ แต่กุญแจยิ่งเยอะยิ่งซับซ้อนขึ้นทุกที จะทำอย่างไรดี? ตัวจัดการรหัสผ่านคือ “ผู้จัดการกุญแจอเนกประสงค์” ของคุณ}

\booksection{หนึ่ง ทำไมคุณถึงต้องมีตัวจัดการรหัสผ่าน?}

คนยุคใหม่มีบัญชีออนไลน์เฉลี่ยตั้งแต่กี่สิบจนถึงเกือบร้อยบัญชี ถ้าทุกบัญชีใช้รหัสผ่านเดียวกัน เว็บไซต์เดียวโดนแฮก บัญชีทั้งหมดก็อันตราย ถ้าแต่ละบัญชีใช้รหัสผ่านซับซ้อนคนละแบบ ก็จำไม่ไหว

ตัวจัดการรหัสผ่านช่วยแก้ทางตันนี้: คุณแค่ต้องจำรหัสผ่านหลักหนึ่งตัว ส่วนที่เหลือมันจำให้คุณ

\booksection{สอง หลักการทำงานของตัวจัดการรหัสผ่าน}

\begin{enumerate}
  \item \textbf{การจัดเก็บแบบเข้ารหัส}: รหัสผ่านทุกตัวถูกเข้ารหัสด้วยความแข็งแรงสูงก่อนเก็บ และมีแค่คุณเท่านั้นที่ถอดได้
  \item \textbf{รหัสผ่านหลัก}: คุณตั้งรหัสผ่านหลักหนึ่งตัวไว้ปลดล็อกคลังรหัสผ่าน ซึ่งตัวรหัสผ่านหลักเองจะไม่ถูกเก็บบนเซิร์ฟเวอร์
  \item \textbf{การกรอกอัตโนมัติ}: ที่หน้าล็อกอินของเว็บหรือแอป ตัวจัดการรหัสผ่านจะกรอกบัญชีและรหัสผ่านให้อัตโนมัติ
  \item \textbf{การสร้างรหัสผ่าน}: สร้างรหัสผ่านแบบสุ่มที่ซับซ้อนให้ได้ เพิ่มความปลอดภัย
\end{enumerate}

\booksection{บทนำเกี่ยวกับหลักการเข้ารหัส}

การทำความเข้าใจหลักการเข้ารหัสของตัวจัดการรหัสผ่าน จะช่วยให้คุณประเมินได้ว่าผู้ให้บริการรายนั้นไว้วางใจได้หรือไม่ ต่อไปนี้จะอธิบายแนวคิดหลักๆ ด้วยภาษาที่เข้าใจง่าย

\booksubsection{การเข้ารหัส AES-256}

AES (Advanced Encryption Standard) คืออัลกอริทึมการเข้ารหัสแบบสมมาตรที่ได้รับการรับรองจากสถาบันมาตรฐานและเทคโนโลยีแห่งชาติสหรัฐอเมริกา (NIST) ซึ่งถูกใช้อย่างแพร่หลายในธนาคาร กองทัพ และหน่วยงานรัฐบาล “256” หมายถึงความยาวของกุญแจที่ 256 บิต — ในทางทฤษฎีแล้วจะต้องใช้เวลามากกว่าอายุของจักรวาลในการเจาะรหัสแบบ Brute Force ตัวจัดการรหัสผ่านยอดนิยม (1Password, Bitwarden, KeePass และอื่นๆ) ล้วนใช้การเข้ารหัส AES-256 กับคลังรหัสผ่านของคุณ

\booksubsection{สถาปัตยกรรมความรู้เป็นศูนย์ (Zero-Knowledge Architecture)}

“ความรู้เป็นศูนย์” ไม่ได้แปลว่าผู้ให้บริการจะไม่เก็บข้อมูลอะไรเลย แต่หมายความว่า\textbf{การเข้ารหัสและถอดรหัสทำบนอุปกรณ์ของคุณเองเท่านั้น เซิร์ฟเวอร์จะเก็บเฉพาะข้อมูลที่เข้ารหัสแล้วเท่านั้น}ขั้นตอนมีดังนี้:

\begin{enumerate}
  \item คุณกรอกรหัสผ่านหลักในไคลเอนต์
  \item ไคลเอนต์จะใช้รหัสผ่านหลักเพื่อสร้างคีย์เข้ารหัส (ขั้นตอนนี้ทำบนมือถือ/คอมพิวเตอร์ของคุณ)
  \item มีเพียงข้อความเข้ารหัสเท่านั้นที่จะถูกอัปโหลดไปยังเซิร์ฟเวอร์
  \item ในครั้งต่อไปที่ล็อกอิน เซิร์ฟเวอร์จะส่งข้อความเข้ารหัสลงมา และไคลเอนต์จะใช้รหัสผ่านหลักถอดรหัสอีกครั้ง
\end{enumerate}

ประเด็นสำคัญ:\textbf{รหัสผ่านหลักจะไม่มีวันออกจากอุปกรณ์ของคุณ}และไม่มีอยู่บนเซิร์ฟเวอร์ด้วย ดังนั้นแม้ผู้ให้บริการจะถูกแฮ็ก สิ่งที่แฮ็กเกอร์ได้ไปก็เพียงข้อความเข้ารหัสที่ถอดไม่ได้ — ยกเว้นว่าพวกเขาจะเดารหัสผ่านหลักของคุณถูก

\begin{tcbinfo}{ทำไม "Zero Knowledge" ถึงหมายความว่าผู้ให้บริการก็มองเห็นรหัสผ่านของคุณไม่ได้?}
เพราะคีย์ถอดรหัสถูกสร้างขึ้นจากรหัสผ่านหลักของคุณ ซึ่งรหัสผ่านหลักนั้นมีอยู่เพียงในความทรงจำของคุณและใช้งานบนอุปกรณ์ของคุณเท่านั้น บนเซิร์ฟเวอร์ของผู้ให้บริการไม่มีทั้งรหัสผ่านหลักและคีย์ถอดรหัส มีเพียงข้อความเข้ารหัสเท่านั้น นี่คือเหตุผลที่ผู้ให้บริการที่ได้มาตรฐานอย่าง Bitwarden, 1Password และอื่นๆ ระบุชัดเจนว่า “พวกเราก็มองเห็นรหัสผ่านของคุณไม่ได้เช่นกัน” — ไม่ใช่ว่าพวกเขาไม่อยากดู แต่ทำไม่ได้ในทางเทคนิค

\end{tcbinfo}

\booksubsection{ฟังก์ชันการได้มาซึ่งคีย์: PBKDF2 และ Argon2}

รหัสผ่านหลักไม่สามารถนำมาใช้เป็นคีย์เข้ารหัสได้โดยตรง — มันจำเป็นต้องผ่านกระบวนการ “ฟังก์ชันการได้มาซึ่งคีย์” หน้าที่ของฟังก์ชันนี้คือ\textbf{ตั้งใจทำให้การคำนวณช้าลง}เพื่อให้ต้นทุนในการเจาะรหัสแบบบรูทฟอร์ซของแฮ็กเกอร์สูงจนเป็นไปไม่ได้

\begin{itemize}
  \item \textbf{PBKDF2}: มาตรฐานเก่าที่ใช้อย่างแพร่หลายใน iCloud คีย์เชน, Bitwarden, LastPass และอื่นๆ โดยใช้การแฮชซ้ำๆ กัน (มักจะมากกว่า 100,000 ครั้ง) เพื่อชะลอการแฮ็ก
  \item \textbf{Argon2}: แชมป์จากการแข่งขันแฮชรหัสผ่านปี 2015 ซึ่งทันสมัยกว่า PBKDF2 นอกจากจะใช้ "พลังการประมวลผล" แล้วยังเพิ่ม "การใช้หน่วยความจำ" ด้วย — ทำให้ฮาร์ดแวร์เฉพาะทางอย่าง GPU/ASIC ต่างก็เร่งความเร็วในการแฮ็กได้ยาก Bitwarden, KeePassXC และอื่นๆ ได้รองรับแล้ว
\end{itemize}

ยิ่งจำนวนรอบการทำซ้ำสูงก็ยิ่งปลอดภัย แต่ความเร็วในการปลดล็อกก็จะช้าลงด้วย ผู้ให้บริการที่ได้มาตรฐานจะเปิดเผยอัลกอริทึมการได้มาซึ่งคีย์และจำนวนรอบการทำซ้ำที่ใช้ เพื่อให้ชุมชนด้านความปลอดภัยตรวจสอบได้

\booksubsection{กรณีศึกษาเตือนใจ: การรั่วไหลของข้อมูล LastPass ในปี 2022}

ในเดือนสิงหาคม 2022 ผู้จัดการรหัสผ่านที่มีชื่อเสียงอย่าง LastPass ได้ประกาศว่าสภาพแวดล้อมการพัฒนาของตนถูกบุกรุก และในเดือนพฤศจิกายนของปีเดียวกันได้เปิดเผยเพิ่มเติมว่า แฮ็กเกอร์ได้ไฟล์สำรองข้อมูลคลังรหัสผ่านของผู้ใช้ (ข้อความเข้ารหัส) ไป บทเรียนสำคัญจากเหตุการณ์นี้คือ:

\begin{itemize}
  \item \textbf{สถาปัตยกรรม Zero Knowledge ได้รักษาเส้นทางไว้}เนื่องจากคลังรหัสผ่านถูกเข้ารหัสเอง แฮ็กเกอร์จึงได้ข้อความเข้ารหัสไปแต่ถอดไม่ได้ — โดยมีเงื่อนไขว่ารหัสผ่านหลักของผู้ใช้ต้องแข็งแรงพอ
  \item \textbf{ผู้ใช้ที่มีรหัสผ่านหลักอ่อนแอเสี่ยงที่สุด}แฮ็กเกอร์สามารถใช้บรูทฟอร์ซออฟไลน์กับรหัสผ่านหลักที่อ่อนแอได้ ดังนั้นผู้ใช้ที่ตั้งรหัสผ่านหลักที่อ่อนแอหรือใช้รหัสผ่านซ้ำกันจึงเผชิญกับความเสี่ยงอย่างแท้จริง
  \item \textbf{ข้อมูลเมตาก็รั่วไหลด้วย}แม้เนื้อหาของรหัสผ่านจะถูกเข้ารหัส แต่ข้อมูลเมตาอย่าง URL ของเว็บไซต์ หรือบันทึกย่อ อาจไม่ได้เข้ารหัส ซึ่งจะเปิดเผยว่าผู้ใช้ใช้บริการใดบ้าง
\end{itemize}

\begin{tcbwarn}{สิ่งที่ได้เรียนรู้จากเหตุการณ์ LastPass}
หนึ่ง รหัสผ่านหลักต้องยาวและไม่ซ้ำใคร (แนะนำให้มากกว่า 16 หลัก ประกอบด้วยตัวพิมพ์เล็ก ตัวพิมพ์ใหญ่ ตัวเลข และสัญลักษณ์); สอง ควรเลือกผู้ให้บริการที่รองรับ Argon2id และมีจำนวนรอบการทำซ้ำสูงเป็นอันดับแรก; สาม สำหรับบัญชีที่มีความสำคัญสูง (ธนาคาร, อีเมล) ต้องเปิดใช้งานการยืนยันตัวตนสองขั้นตอนแยกต่างหาก อย่าพึ่งพาผู้จัดการรหัสผ่านเพียงอย่างเดียว

\end{tcbwarn}

\booksection{แนะนำผู้จัดการรหัสผ่านยอดนิยม}

\booksubsection{มากับระบบ:}

\begin{itemize}
  \item \textbf{iCloud Keychain (พวงกุญแจ iCloud)}: ทางเลือกแรกของผู้ใช้ระบบนิเวศ Apple เชื่อมรวมกับ Safari และแอปต่าง ๆ ได้ไร้รอยต่อ ฟรีและปลอดภัย ข้อเสียคือใช้ข้ามแพลตฟอร์มไม่สะดวก
  \item \textbf{Google Password Manager}: สะดวกสำหรับผู้ใช้ Android เชื่อมรวมกับ Chrome และระบบ Android ใช้ได้ฟรี
  \item \textbf{Huawei Password Vault (ตู้นิรภัยรหัสผ่าน Huawei)}: มากับระบบ HarmonyOS เหมาะกับผู้ใช้ Huawei
\end{itemize}

\booksubsection{เครื่องมือระดับมืออาชีพจากบุคคลที่สาม:}

\begin{itemize}
  \item \textbf{Bitwarden}: โอเพนซอร์สฟรี ฟังก์ชันแข็งแรงครบครัน รองรับทุกแพลตฟอร์ม แนะนำอย่างยิ่งสำหรับผู้ใช้ที่ใส่ใจความเป็นส่วนตัว
  \item \textbf{1Password}: อินเทอร์เฟซสวยงาม ฟังก์ชันจัดเต็ม เหมาะกับครอบครัวและทีม เสียเงิน
  \item \textbf{LastPass}: ตัวจัดการรหัสผ่านรุ่นเก่า เวอร์ชันฟรีมีฟังก์ชันจำกัด
  \item \textbf{KeePass}: เก็บข้อมูลในเครื่องล้วน ๆ โอเพนซอร์สฟรี เหมาะกับสายเกจ
\end{itemize}

\booksection{การประเมินความแข็งแรงของรหัสผ่านอย่างแท้จริง}

ผู้จัดการรหัสผ่านมักจะแสดงแถบบ่งชี้ “ความแข็งแรงของรหัสผ่าน” (แดง/เหลือง/เขียว) แต่การประเมินนี้มักจะดูเพียงความยาวและประเภทของตัวอักษรเท่านั้น ซึ่งไม่ครอบคลุม การตัดสินว่ารหัสผ่านปลอดภัยจริงหรือไม่ จำเป็นต้องพิจารณา “เอนโทรปี” และ “ว่าเคยถูกรั่วไหลหรือไม่”

\booksubsection{เอนโทรปี (Entropy): การวัดความแข็งแรงของรหัสผ่านทางวิทยาศาสตร์}

“เอนโทรปี” ใช้วัดความไม่แน่นอนของรหัสผ่าน โดยมีหน่วยเป็นบิต ยิ่งเอนโทรปีสูง จำนวนครั้งในการลองผิดทดลองเพื่อบรูทฟอร์ซก็ยิ่งมากขึ้น

\begin{itemize}
  \item \textbf{"123456"}: เอนโทรปีต่ำมาก แฮ็กได้ในพริบตา
  \item \textbf{"P@ssw0rd"}: ดูเหมือนจะซับซ้อน แต่กลับปรากฏอยู่ในทุกรายการของการโจมตีแบบพจนานุกรม เอนโทรปีจึงยังต่ำอยู่
  \item \textbf{"correct horse battery staple"}: คำศัพท์สุ่ม 4 คำ มีเอนโทรปีประมาณ 44 บิต ความยากในการแฮ็กสูงกว่า "P@ssw0rd" อย่างมาก
  \item \textbf{อักขระสุ่ม 16 หลัก (ประกอบด้วยตัวพิมพ์เล็ก+ตัวพิมพ์ใหญ่+ตัวเลข+สัญลักษณ์)}: มีเอนโทรปีประมาณ 100 บิต ด้วยกำลังการประมวลผลในปัจจุบันไม่สามารถบรูทฟอร์ซได้
\end{itemize}

บทสรุป:\textbf{“ความยาว + ความสุ่ม” สำคัญกว่า “สัญลักษณ์ที่ซับซ้อน”}นี่คือเหตุผลว่าทำไมรหัสผ่านที่ยาวและสุ่มซึ่งสร้างโดยผู้จัดการรหัสผ่านถึงปลอดภัยอย่างแท้จริง

\booksubsection{HaveIBeenPwned: ตรวจสอบว่ารหัสผ่านของคุณถูกรั่วไหลหรือไม่}

ฐานข้อมูลที่ดูแลโดยนักวิจัยด้านความปลอดภัย Troy Hunt\href{https://haveibeenpwned.com}{haveibeenpwned.com}รวบรวมบัญชีและรหัสผ่านที่รั่วไหลจากการละเมิดข้อมูลครั้งใหญ่ในอดีต คุณสามารถ:

\begin{enumerate}
  \item กรอกอีเมลเพื่อตรวจสอบว่าบัญชีของคุณเคยถูกเปิดเผยในการรั่วไหลครั้งใดหรือไม่
  \item ใช้ฟังก์ชัน Pwned Passwords เพื่อค้นหาว่ารหัสผ่านใดรหัสผ่านหนึ่งอยู่ในฐานข้อมูลที่รั่วไหลที่ทราบหรือไม่ (ผู้จัดการรหัสผ่านและเบราว์เซอร์ได้ฝังการค้นหานี้ไว้ในตัวแล้ว)
\end{enumerate}

หากรหัสผ่านปรากฏอยู่ในฐานข้อมูลที่รั่วไหล\textbf{ไม่ว่ามันจะดูซับซ้อนแค่ไหนก็ต้องเปลี่ยนทันที}เพราะการโจมตีแบบพจนานุกรมจะลองรหัสผ่านที่ทราบเหล่านี้เป็นอันดับแรก

\begin{tcbinfo}{หลักการค้นหาแบบ k-Anonymity}
กังวลว่าการกรอกรหัสผ่านใน HaveIBeenPwned จะรั่วไหลอีกครั้ง? ไม่ต้องกังวล การค้นหาสมัยใหม่ใช้โปรโตคอล “k-Anonymity”: เบราว์เซอร์จะส่งเพียง 5 หลักแรกของแฮชรหัสผ่านไปยังเซิร์ฟเวอร์ จากนั้นเซิร์ฟเวอร์จะส่งรายการแฮชที่ขึ้นต้นด้วย 5 หลักนั้นกลับมาทั้งหมด แล้วจึงเปรียบเทียบบนเครื่องของคุณ กล่าวคือ เซิร์ฟเวอร์จะไม่มีวันเห็นแฮชรหัสผ่านฉบับเต็มของคุณ การตรวจจับรหัสผ่านที่รั่วไหลของ 1Password และ Bitwarden ก็ทำงานแบบนี้

\end{tcbinfo}

\booksubsection{ฟังก์ชันการตรวจสอบที่ฝังอยู่ในผู้จัดการรหัสผ่าน}

\begin{itemize}
  \item \textbf{1Password Watchtower}ตรวจสอบแบบเรียลไทม์สำหรับรหัสผ่านที่อ่อนแอ รหัสผ่านที่ซ้ำกัน รหัสผ่านที่รั่วไหล รหัสผ่านที่หมดอายุ และเว็บไซต์ที่ไม่ได้ใช้ 2FA พร้อมแสดงผลแบบจัดระดับสีแดง/ส้ม/เหลือง
  \item \textbf{Bitwarden Vault Health}:เวอร์ชันฟรีสามารถดูรหัสผ่านที่อ่อนแอและรหัสผ่านที่ซ้ำกันได้ เวอร์ชันพรีเมียมสามารถตรวจสอบรหัสผ่านที่รั่วไหลออกสู่สาธารณะได้ (อ้างอิงจาก HaveIBeenPwned)
  \item \textbf{"คำแนะนำด้านความปลอดภัย" ของ iCloud คีย์เชน}:iOS ไปที่ การตั้งค่า รหัสผ่าน คำแนะนำด้านความปลอดภัย ระบบจะแจ้งเตือนว่ารหัสผ่านใดปรากฏในการรั่วไหลที่ทราบ และรหัสผ่านใดที่มีการใช้ซ้ำ
  \item \textbf{"เครื่องมือตรวจสอบรหัสผ่าน" ของ Google รหัสผ่าน}:ภายในตัว Chrome สามารถตรวจสอบรหัสผ่านที่รั่วไหล อ่อนแอ และซ้ำกันได้ในคลิกเดียว
\end{itemize}

\begin{tcbpractice}{ตรวจสอบตอนนี้: รหัสผ่านของคุณปลอดภัยหรือไม่}
\begin{enumerate}
  \item เปิด haveibeenpwned.com แล้วกรอกอีเมลที่ใช้บ่อย เพื่อดูว่าปรากฏในเหตุการณ์การรั่วไหลใดบ้าง
  \item ผู้ใช้ iPhone: ไปที่ การตั้งค่า รหัสผ่าน คำแนะนำด้านความปลอดภัย แล้วจัดการคำเตือนสีแดงทีละรายการ
  \item ผู้ใช้ Chrome: ไปที่ การตั้งค่า รหัสผ่าน ตรวจสอบรหัสผ่าน แล้วแก้ไขตามคำแนะนำ
  \item ผู้ใช้ 1Password/Bitwarden: เปิด Watchtower/Health Reports โดยจัดลำดับเรื่อง "รั่วไหล" และ "ซ้ำกัน" ก่อน
\end{enumerate}

\end{tcbpractice}

\booksection{Passkey: อนาคตของการเข้าสู่ระบบแบบไร้รหัสผ่าน}

ในช่วงปี 2024-2026 เทคโนโลยีรหัสผ่านกำลังเผชิญกับการปฏิวัติครั้งใหญ่ —\textbf{Passkey (คีย์การเข้าถึง)}คาดว่าจะสามารถแทนที่รหัสผ่านแบบดั้งเดิมได้อย่างสิ้นเชิง โดยมี Apple, Google, และ Microsoft สามแพลตฟอร์มหลักร่วมกันผลักดัน อีกทั้งยังได้รับการสนับสนุนจากพันธมิตร FIDO และ World Wide Web Consortium (W3C) นี่ถือเป็นการเปลี่ยนแปลงวิธีการเข้าสู่ระบบครั้งใหญ่ที่สุดในรอบหลายสิบปี

\booksubsection{Passkey คืออะไร?}

Passkey อิงตามมาตรฐาน FIDO2 / WebAuthn โดยใช้\textbf{คู่กุญแจบนอุปกรณ์}แทนรหัสผ่าน หลักการทำงานสามารถอธิบายสั้นๆ ได้ดังนี้:

\begin{itemize}
  \item โทรศัพท์ของคุณจะสร้างคู่กุญแจขึ้นมา 1 ชุด:\textbf{กุญแจสาธารณะ}(ลงทะเบียนกับเว็บไซต์) +\textbf{กุญแจส่วนตัว}(จะถูกเก็บไว้ในชิปความปลอดภัยของอุปกรณ์คุณตลอดไป)
  \item เมื่อเข้าสู่ระบบ เว็บไซต์จะส่ง "คำท้าทาย" มาให้ แล้วโทรศัพท์ของคุณจะใช้กุญแจส่วนตัวเซ็นกำกับ จากนั้นเว็บไซต์จะใช้กุญแจสาธารณะในการตรวจสอบ
  \item ก่อนการเซ็นกำกับ จะต้องมีการยืนยันตัวตนผ่านชีวภาพ (Face ID / ลายนิ้วมือ) หรือ PIN ของอุปกรณ์
  \item ตลอดกระบวนการนี้,\textbf{ไม่มีรหัสผ่านใดๆ ถูกส่งผ่านเครือข่าย}
\end{itemize}

\booksubsection{Passkey เทียบกับรหัสผ่านแบบดั้งเดิม}

\begin{center}
\begin{tabularx}{\linewidth}{|X|X|X|}
\hline
\textbf{มิติข้อมูล} & \textbf{รหัสผ่านแบบดั้งเดิม} & \textbf{Passkey} \tabularnewline
\hline
ป้องกันฟิชชิง & แย่ เว็บไซต์ฟิชชิงสามารถขโมยรหัสผ่านได้ & ป้องกันฟิชชิงได้โดยธรรมชาติ กุญแจส่วนตัวจะเซ็นกำกับเฉพาะโดเมนจริงเท่านั้น \tabularnewline
ผลกระทบจากการรั่วไหลของเซิร์ฟเวอร์ & แฮชรหัสผ่านอาจถูกแกะรอยได้ & เซิร์ฟเวอร์เก็บเพียงกุญแจสาธารณะ การรั่วไหลไม่มีผลกระทบ \tabularnewline
ความสะดวกสบาย & ต้องจำ/พิมพ์/จัดการ & เข้าสู่ระบบใน 1 วินาทีด้วย Face ID/ลายนิ้วมือ \tabularnewline
โจมตีซ้ำ (Replay Attack) & อาจถูกดักจับและเล่นซ้ำโดยคนกลาง & ลายเซ็นการเข้าสู่ระบบแต่ละครั้งใช้ครั้งเดียว ไม่สามารถนำไปเล่นซ้ำได้ \tabularnewline
\hline
\end{tabularx}
\end{center}

\booksubsection{การรองรับในแต่ละแพลตฟอร์ม}

\begin{itemize}
  \item \textbf{Passkey ของ Apple iCloud คีย์เชน}:รองรับในตัวตั้งแต่ iOS 16+ / macOS Ventura+ Passkey สามารถซิงค์ระหว่างอุปกรณ์ที่เข้าสู่ระบบ Apple ID เดียวกันทั้งหมด เมื่อใช้งานบนอุปกรณ์เครื่องใหม่ ไม่ต้องสร้างใหม่
  \item \textbf{Passkey ของ Google รหัสผ่าน}:รองรับในตัวบน Android + Chrome Passkey จะถูกเก็บไว้ในบัญชี Google และซิงค์ข้ามอุปกรณ์ Android ตั้งแต่ปี 2023 เป็นต้นมา บัญชีส่วนตัวของ Google รองรับการเข้าสู่ระบบด้วย Passkey แล้ว
  \item \textbf{Microsoft}:Windows 11 / Edge รองรับ Passkey แล้ว บริการต่างๆ ของ Microsoft เช่น บัญชี Microsoft, Outlook, Office ก็ค่อยๆ เปิดให้ใช้งานการเข้าสู่ระบบด้วย Passkey และสามารถจัดการได้ใน Microsoft Authenticator
  \item \textbf{1Password / Bitwarden}:รองรับการสร้าง จัดเก็บ และซิงค์ข้ามแพลตฟอร์มของ Passkey ทำให้คุณไม่ถูกจำกัดอยู่แค่ในระบบนิเวศเดียว
\end{itemize}

\begin{tcbinfo}{การซิงค์ข้ามอุปกรณ์ถือเป็นความก้าวหน้าที่สำคัญที่สุด}
จุดอ่อนที่ใหญ่ที่สุดของ Passkey ในช่วงแรกคือ “ผูกกับอุปกรณ์เครื่องเดียว” — เปลี่ยนมือถือก็ต้องสร้างใหม่ ตั้งแต่ปี 2023 Apple, Google และ Microsoft ได้ทำการซิงค์ Passkey ผ่านคลาวด์: iCloud ซิงค์ Keychain ข้ามอุปกรณ์ Apple, Google Password Manager ซิงค์ข้ามอุปกรณ์ Android ส่วนบริการของบุคคลที่สามอย่าง 1Password ก็เชื่อมต่อการซิงค์ข้ามระบบนิเวศได้แล้ว นั่นหมายความว่า Passkey มีความพร้อมในการใช้งานจริงเพื่อทดแทนรหัสผ่านแล้ว

\end{tcbinfo}

\booksubsection{วิธีสร้างและใช้ Passkey บนมือถือ}

\textbf{ขั้นตอนบน iOS (iPhone):}

\begin{enumerate}
  \item ไปที่หน้าความปลอดภัยของบัญชีเว็บไซต์ที่รองรับ Passkey (เช่น Google, Microsoft, Amazon, eBay, PayPal) แล้วเลือก "เพิ่ม Passkey / Add passkey"
  \item ระบบจะแสดงหน้าต่างขอ Face ID / Touch ID เพื่อยืนยันตัวตน เมื่อผ่านแล้ว Passkey จะถูกสร้างและบันทึกลงใน iCloud Keychain โดยอัตโนมัติ
  \item ครั้งต่อไปที่เข้าสู่ระบบเว็บไซต์ดังกล่าว หลังกรอกชื่อผู้ใช้ให้เลือก "ใช้ iPhone เข้าสู่ระบบ" แค่สแกน Face ID ก็เสร็จ
  \item ทางเข้าจัดการ: การตั้งค่า รหัสผ่าน ค้นหาเว็บไซต์ที่ต้องการ จะเห็นรายการ Passkey
\end{enumerate}

\textbf{ขั้นตอนบน Android:}

\begin{enumerate}
  \item ไปที่หน้าความปลอดภัยของบัญชีเว็บไซต์ที่รองรับ Passkey แล้วเลือก "เพิ่ม Passkey"
  \item Google Password Manager จะเด้งขึ้นมา เลือกบัญชี Google ของคุณ ยืนยันด้วยการล็อกหน้าจอ (ลายนิ้วมือ/รูปแบบ/PIN)
  \item Passkey จะถูกบันทึกไว้ใน Google Password Manager โดยอัตโนมัติ และซิงค์ข้ามอุปกรณ์ Android ที่ใช้บัญชีเดียวกัน
  \item ครั้งต่อไปที่เข้าสู่ระบบ ให้เลือก "ใช้ Passkey" แค่ปลดล็อกหน้าจอก็เสร็จ
  \item ทางเข้าจัดการ: การตั้งค่า Android รหัสผ่าน รหัสผ่านคีย์ และการกรอกอัตโนมัติ ตัวจัดการรหัสผ่าน google.com
\end{enumerate}

\booksubsection{การเข้าสู่ระบบ Passkey ข้ามอุปกรณ์}

หากคุณเข้าสู่ระบบเว็บไซต์บนคอมพิวเตอร์ Windows ที่มี Passkey อยู่บน iPhone เพียงเครื่องเดียว สามารถใช้การสแกน QR Code เพื่ออนุญาตได้:

\begin{enumerate}
  \item ที่หน้าเข้าสู่ระบบของเว็บไซต์ ให้เลือก "ใช้ Passkey"
  \item เลือก "ใช้อุปกรณ์อื่น"
  \item คอมพิวเตอร์จะแสดง QR Code ให้ใช้ iPhone สแกน
  \item iPhone จะเด้งขึ้นมาขอยืนยัน Face ID เมื่อยืนยันแล้วเว็บไซต์จะเข้าสู่ระบบบนคอมพิวเตอร์สำเร็จ
\end{enumerate}

กระบวนการนี้ใช้การเชื่อมต่อแบบเข้ารหัสลับสองชั้นผ่าน Bluetooth และอินเทอร์เน็ต โดย Private Key ของ Passkey จะไม่ออกจาก iPhone ของคุณ

\begin{tcbpractice}{ลองเลยตอนนี้: สัมผัสประสบการณ์การเข้าสู่ระบบไร้รหัสผ่าน}
\begin{enumerate}
  \item เข้าสู่ระบบบัญชี Google ความปลอดภัย Passkey เพิ่ม Passkey เพื่อทดลองเข้าสู่ระบบด้วย Face ID/ลายนิ้วมือ
  \item เข้าสู่ระบบบัญชี Microsoft ความปลอดภัย ตัวเลือกความปลอดภัยขั้นสูง เพิ่ม Passkey
  \item ลองใช้บนบริการที่รองรับ Passkey แล้ว เช่น Taobao, JD, eBay, PayPal
  \item เก็บรหัสผ่านเดิมไว้เป็นตัวสำรอง แล้วค่อย ๆ ย้ายบัญชีที่มีความสำคัญสูงไปใช้ Passkey
\end{enumerate}

\end{tcbpractice}

\begin{tcbwarn}{คำแนะนำช่วงการเปลี่ยนผ่าน}
Passkey ยังไม่ได้ทดแทนรหัสผ่านแบบเต็มรูปแบบ — เว็บไซต์หลายแห่งยังอยู่ในช่วง “Passkey + รหัสผ่าน” ที่มีให้ใช้พร้อมกัน ขอแนะนำ: เปิดใช้งาน Passkey สำหรับบัญชีที่มีความสำคัญสูงก่อน แต่\textbf{อย่าลบรหัสผ่านเดิมทิ้งทันที}ให้เก็บไว้เป็นตัวสำรอง พร้อมทั้งยังคงใช้ Password Manager ต่อไป เพราะเว็บไซต์ไม่ได้รองรับ Passkey ทั้งหมด ให้มอง Passkey เป็น “สิ่งที่เสริมความแข็งแกร่งให้ Password Manager” มากกว่า “การทดแทนทันที”

\end{tcbwarn}

\booksection{บทที่ 7 · ฮาร์ดแวร์ความปลอดภัย (Hardware Security Key)}

หากคุณมีความต้องการด้านความปลอดภัยสูงกว่าผู้ใช้ทั่วไป — เช่น ถือครองสินทรัพย์ดิจิทัลจำนวนมาก เป็นผู้บริหาร/นักข่าว/นักเคลื่อนไหว หรือต้องปกป้องบัญชีหลักขององค์กร — ฮาร์ดแวร์ความปลอดภัยถือเป็นทางเลือกที่น่าเชื่อถือกว่ารหัสยืนยันทาน SMS

\booksubsection{ฮาร์ดแวร์ความปลอดภัยคืออะไร?}

ฮาร์ดแวร์คีย์คืออุปกรณ์ทางกายภาพ (ลักษณะคล้าย USB Flash Drive) โดยมีชิปความปลอดภัยภายในเก็บ Private Key ของคุณไว้ ซึ่งต้องสัมผัสทางกายภาพหรือใช้ NFC จึงจะยืนยันตัวตนได้ ผลิตภัณฑ์ที่เป็นตัวแทน:

\begin{itemize}
  \item \textbf{YubiKey}(บริษัท Yubico): เป็นที่นิยมที่สุด รองรับ FIDO2/WebAuthn, U2F, OTP, PIV, OpenPGP มีหลายอินเทอร์เฟซเช่น USB-C / USB-A / NFC รุ่น Series 5 ราคา 50-70 ดอลลาร์
  \item \textbf{Google Titan}: ผลิตโดย Google เอง รองรับ USB-C + NFC / Bluetooth ราคาประมาณ 30-40 ดอลลาร์
  \item \textbf{OnlyKey}: ทางเลือกโอเพนซอร์ส รองรับหลายโปรโตคอล
  \item \textbf{ทางเลือกจากผู้ผลิตในประเทศจีน}: เช่น ePasskey ของ Feitian ราคาประหยัดกว่า
\end{itemize}

\booksubsection{ทำไมฮาร์ดแวร์คีย์จึงปลอดภัยกว่ารหัสยืนยันทาน SMS?}

\begin{itemize}
  \item \textbf{ป้องกันฟิชชิง}: เช่นเดียวกับ Passkey ฮาร์ดแวร์คีย์จะตอบสนองเฉพาะโดเมนจริงเท่านั้น เว็บไซต์หลอกลวงไม่สามารถดึงรหัสยืนยันได้
  \item \textbf{ต้านทานการปล้น SIM}: รหัสยืนยันทาน SMS อาจถูกดักจับโดยการโจมตีแบบ SIM-swap ได้ แต่ฮาร์ดแวร์คีย์ไม่ต้องใช้หมายเลขโทรศัพท์
  \item \textbf{ต้านทาน Man-in-the-Middle}: โปรโตคอล FIDO2 จะมี Challenge Value แทรกในการเซ็นชื่อทุกครั้ง ทำให้ไม่สามารถนำไปใช้ซ้ำได้
  \item \textbf{การแยกทางกายภาพ}: Private Key จะอยู่ในชิปฮาร์ดแวร์เสมอ แม้มัลแวร์จะเจาะระบบคอมพิวเตอร์ของคุณได้ก็ไม่สามารถอ่านคีย์นั้นได้
\end{itemize}

\booksubsection{สถานการณ์ที่เหมาะสม}

\begin{itemize}
  \item บัญชีการแลกเปลี่ยนสินทรัพย์ดิจิทัล การกู้คืนกระเป๋า Wallet
  \item SSO ขององค์กร, VPN สำหรับทำงานทางไกล, ล็อกอินเซิร์ฟเวอร์ผ่าน SSH
  \item อีเมล, GitHub, โซเชียลมีเดีย และ "บัญชีหลัก" อื่น ๆ — หากถูกบุกรุก จะสามารถรีเซ็ตบัญชีอื่น ๆ ทั้งหมดได้
  \item นักข่าว ทนายความ ผู้บริหารระดับสูง และบุคคลที่มีความอ่อนไหวสูง
\end{itemize}

\booksubsection{วิธีตั้งค่า Hardware Key เป็นการยืนยันตัวตน 2 ขั้นตอน}

ยกตัวอย่างบัญชี Google + YubiKey:

\begin{enumerate}
  \item ซื้อ YubiKey (แนะนำให้ซื้อครั้งละ 2 อัน หนึ่งอันไว้ใช้งานประจำ อีกอันไว้สำรอง โดยเก็บไว้คนละสถานที่)
  \item เข้าสู่ระบบ myaccount.google.com ความปลอดภัย การยืนยันตัวตน 2 ขั้นตอน
  \item เลือก "เพิ่มคีย์ความปลอดภัย" เสียบ YubiKey เข้ากับพอร์ต USB-C หรือนำไปแตะใกล้พื้นที่เซ็นเซอร์ NFC ของโทรศัพท์
  \item ทำตามคำแนะนำโดยสัมผัสที่จุดสัมผัสโลหะของ YubiKey เพื่อทำการลงทะเบียนให้เสร็จสมบูรณ์
  \item แนะนำให้ลงทะเบียนคีย์สำรองตัวที่สองด้วย
  \item ในการตั้งค่า "การยืนยันตัวตน 2 ขั้นตอน" ให้ตั้งค่า Hardware Key เป็นวิธียืนยันหลัก และใช้รหัส SMS เป็นตัวสำรอง
\end{enumerate}

เว็บไซต์ที่รองรับหลัก: Google, GitHub, Microsoft, Apple ID, Facebook, Twitter, Dropbox, GitLab, AWS, Coinbase, Binance และอื่นๆ

\begin{tcbinfo}{กลยุทธ์คู่คีย์}
ความเสี่ยงที่ใหญ่ที่สุดของ Hardware Key คือ “การสูญหาย” ขอแนะนำ: (1) ลงทะเบียนคีย์ที่แตกต่างกัน 2 อันพร้อมกัน; (2) หนึ่งอันพกติดตัว (คล้องกับพวงกุญแจ), อีกอันเก็บไว้ในตู้เซฟที่บ้านหรือตู้นิรภัยของธนาคาร; (3) เว็บไซต์บางแห่งอนุญาตให้พิมพ์ “รหัสกู้คืนสำรอง” ออกมาเก็บแบบออฟไลน์ได้ วิธีนี้จะทำให้แม้คีย์หลักหาย คุณก็จะไม่ถูกล็อกออกจากบัญชีอย่างถาวร

\end{tcbinfo}

\begin{tcbwarn}{ระมัดระวังในการเลือกช่องทางซื้อ}
ต้องซื้อจากร้านค้าทางการหรือตัวแทนจำหน่ายที่ได้รับอนุญาตเท่านั้น ห้ามซื้อคีย์มือสองหรือที่มาไม่ชัดเจน—ตามทฤษฎีแล้ว ฮาร์ดแวร์ที่ถูกตั้งค่ามาก่อนอาจถูกแทรกช่องโหว่แบ็กดอร์ไว้แล้ว เฟิร์มแวร์ของ YubiKey ไม่สามารถอัปเดตได้หลังจากออกจากโรงงาน แต่ Yubico จะส่งประกาศความปลอดภัยอย่างต่อเนื่อง เพียงแค่ติดตามเว็บไซต์ทางการของพวกเขาก็พอ

\end{tcbwarn}

\booksection{การแชร์รหัสผ่านในครอบครัว/ทีม}

ในสถานการณ์ครอบครัว Netflix, Wi-Fi ในบ้าน, หน้า backend ของเราเตอร์, รหัสผ่านกลอนอัจฉริยะ จำเป็นต้องใช้ร่วมกันทั้งครอบครัว; ในสถานการณ์ทีม เซิร์ฟเวอร์, ฐานข้อมูล, บัญชีโซเชียลมีเดีย จำเป็นต้องให้หลายคนร่วมกันทำงาน\textbf{การส่งรหัสผ่านผ่าน WeChat หรือเขียนไว้ในเอกสารที่แชร์กันถือเป็นวิธีที่เสี่ยงสูง}—เพราะเมื่อประวัติการแชทหรือสิทธิ์การเข้าถึงเอกสารรั่วไหล รหัสผ่านก็จะรั่วไหลไปด้วย

\booksubsection{“ตู้นิรภัยร่วม” (Vault) ของ 1Password}

1Password คือผลิตภัณฑ์มาตรฐานสำหรับการแชร์ในครอบครัวและทีม กลไกหลักคือ\textbf{ตู้นิรภัยหลายใบ}ก่อนพูดบทสรุป:

\begin{itemize}
  \item \textbf{ตู้นิรภัยส่วนตัว}:มีแค่คุณที่เห็น เช่น บัญชีส่วนตัวอย่างธนาคาร อีเมล
  \item \textbf{ตู้นิรภัยร่วม}:สมาชิกในครอบครัวหรือทีมดูได้ทุกคน เช่น Netflix, Wi-Fi ในบ้าน, บัตรเครดิตร่วม
  \item \textbf{การแบ่งระดับสิทธิ์}:เวอร์ชันทีมสามารถตั้งค่าให้สมาชิกบางคน "ดูได้อย่างเดียว" หรือ "ดูได้แต่แก้ไขไม่ได้"
\end{itemize}

สมาชิกในครอบครัวหรือทีมไม่จำเป็นต้องแชร์ Master Password แต่ละคนใช้บัญชีของตัวเองเข้าสู่ระบบ แต่สามารถดูเนื้อหาในตู้นิรภัยร่วมที่ได้รับอนุญาตได้ เมื่อรหัสผ่านเปลี่ยนแปลง สมาชิกทุกคนจะซิงค์แบบเรียลไทม์

\booksubsection{ฟังก์ชัน “องค์กร” ของ Bitwarden}

เวอร์ชันฟรีของ Bitwarden รองรับการแชร์รหัสผ่านส่วนบุคคล (แชร์รายการเดียวแบบจำกัดเวลาผ่านฟังก์ชัน “Send”); เวอร์ชันเสียเงินสามารถสร้าง “องค์กร” เพื่อแชร์รายการต่างๆ ให้กับสมาชิกในครอบครัวหรือเพื่อนร่วมทีมได้ โดยแบ่งสิทธิ์ได้ละเอียด Bitwarden ยังมีตัวเลือก Self-hosted ซึ่งเหมาะสำหรับองค์กรที่มีความต้องการเรื่องอธิปไตยของข้อมูลสูง

\booksubsection{สถานการณ์ครอบครัว: วิธีปลอดภัยในการแชร์รหัสผ่าน Netflix, Wi-Fi และอื่นๆ}

\begin{itemize}
  \item \textbf{Netflix / เว็บไซต์วิดีโอ}:เก็บรหัสผ่านเข้าสู่ระบบไว้ใน "ตู้นิรภัยร่วม" ของ 1Password เวอร์ชันครอบครัว สมาชิกแต่ละคนเข้าสู่ระบบด้วยตัวเองได้เลย ไม่ต้องส่งกันในกลุ่ม
  \item \textbf{รหัสผ่าน Wi-Fi ในบ้าน}:อุปกรณ์ iOS สามารถแชร์รหัสผ่าน Wi-Fi ให้แขกที่มาเยือนได้โดยตรงผ่าน AirDrop โดยไม่ต้องพูดออกมาเป็นเสียง; Android 12+ ก็มีฟังก์ชัน "แชร์ Wi-Fi ใกล้เคียง" เช่นกัน
  \item \textbf{กลอนอัจฉริยะ, เครื่องใช้ไฟฟ้าอัจฉริยะ}:พยายามใช้ฟังก์ชัน "สมาชิกในครอบครัว" ในแอปของผู้ผลิตเพื่อสร้างบัญชีแยกให้แต่ละคน แทนที่จะแชร์บัญชีผู้ดูแลระบบร่วมกันอันเดียว
  \item \textbf{หน้า backend ของเราเตอร์ในบ้าน}:เก็บรหัสผ่านผู้ดูแลระบบไว้ในตู้นิรภัยร่วมของ 1Password และเปิด SSID แยกต่างหากสำหรับ Wi-Fi แขก
\end{itemize}

\begin{tcbinfo}{ความคุ้มค่าของ 1Password เวอร์ชันครอบครัว}
1Password Families ราคาประมาณ 4.99 ดอลลาร์สหรัฐ/เดือน รองรับสมาชิกในครอบครัวสูงสุด 5 คน แต่ละคนจะมีตู้นิรภัยส่วนตัวอิสระ + ตู้นิรภัยร่วม 4 ใบ คิดเป็นค่าเฉลี่ยต่อคนไม่ถึง 1 ดอลลาร์สหรัฐ/เดือน ถือเป็นโซลูชันที่ดีที่สุดสำหรับการแชร์ในครอบครัวในขณะนี้ เวอร์ชันครอบครัวของ Bitwarden ราคา 10 ดอลลาร์สหรัฐ/ปี รองรับ 6 คน คุ้มค่ากว่า แต่ประสบการณ์การใช้งานอินเทอร์เฟซอาจจะด้อยกว่าเล็กน้อย

\end{tcbinfo}

\booksubsection{สถานการณ์ทีม: การแชร์รหัสผ่านเซิร์ฟเวอร์/ฐานข้อมูล}

\begin{itemize}
  \item \textbf{1Password Teams / Business}:สร้างตู้นิรภัยหลายใบตามแผนกหรือทีมโปรเจกต์ ควบคุมสิทธิ์แบบละเอียด และสามารถตรวจสอบพฤติกรรมการเข้าถึงทั้งหมดได้
  \item \textbf{Bitwarden Teams / Enterprise}:โซลูชันโอเพนซอร์ส สามารถ Self-hosted ได้ เวอร์ชันองค์กรรองรับการเชื่อมต่อ SSO และการซิงค์ไดเรกทอรี
  \item \textbf{หลักการแยกหน้าที่}:รหัสผ่านผู้ดูแลระบบไม่ควรอยู่กับคนเพียงคนเดียว—ควรตั้ง "ผู้ติดต่อฉุกเฉิน" ไว้ 2-3 คน เพื่อให้แน่ใจว่าเมื่อบุคลากรหลักลาออกหรือพักงาน ธุรกิจจะไม่หยุดชะงัก
  \item \textbf{กระบวนการเมื่อลาออก}:เมื่อพนักงานลาออก ให้เพิกถอนสิทธิ์การเข้าถึง Password Manager ทันที และหากจำเป็นต้องเปลี่ยนรหัสผ่านที่มีความอ่อนไหวสูงทั้งหมดที่พนักงานคนนั้นเคยสัมผัส
\end{itemize}

\booksubsection{การเข้าถึงฉุกเฉิน (Emergency Access): หากเกิดเหตุไม่คาดฝันกับคุณ}

ลองนึกภาพสถานการณ์นี้: คุณเป็นคนเดียวที่จัดการบัญชีสำคัญทั้งหมดของครอบครัว แต่หากคุณเข้าโรงพยาบาลกะทันหันหรือเสียชีวิตลง สมาชิกในครอบครัวจะไม่สามารถเข้าสู่ระบบบัญชีธนาคาร อีเมล หรือบัญชีประกันภัยได้\textbf{ฟีเจอร์ “การเข้าถึงฉุกเฉิน”}ถูกออกแบบมาเพื่อเรื่องนี้โดยเฉพาะ — อนุญาตให้คนในครอบครัว/เพื่อนที่ไว้วางใจ เข้าถึงตู้นิรภัยรหัสผ่านของคุณเมื่อคุณไม่สามารถตอบสนองได้

\textbf{การตั้งค่าการเข้าถึงฉุกเฉินของ Bitwarden:}

\begin{enumerate}
  \item เข้าสู่ระบบ Bitwarden Web Vault การตั้งค่า การเข้าถึงฉุกเฉิน
  \item เพิ่มอีเมลของผู้ติดต่อที่เชื่อถือได้ แล้วเชิญให้แก่ยอมรับ
  \item กำหนด "ระยะเวลารอ" (แนะนำ 7-14 วัน): ช่วงเวลาตั้งแต่อีกฝ่ายส่งคำขอจนถึงเวลาที่คุณยืนยัน
  \item หากคุณไม่ตอบสนองในช่วงเวลารอ (ปฏิเสธหรือไม่ดำเนินการใดๆ) เมื่อหมดเวลาแล้ว อีกฝ่ายจะได้รับสิทธิ์การเข้าถึงโดยอัตโนมัติ
  \item คุณสามารถปฏิเสธคำขอหรือเพิกถอนสิทธิ์การเข้าถึงฉุกเฉินได้ตลอดเวลา
\end{enumerate}

\textbf{การเข้าถึงฉุกเฉินของ 1Password:}

\begin{enumerate}
  \item 1Password ไม่ได้เรียกว่า "การเข้าถึงฉุกเฉิน" โดยตรง แต่ทำได้ผ่านฟีเจอร์ "องค์กรครอบครัว/ทีม" — หลังจากเพิ่มสมาชิกในครอบครัวเข้าสู่แบบครอบครัวแล้ว ให้นำรายการบางส่วนใส่ใน "ตู้นิรภัยที่แชร์กัน"
  \item สำหรับผู้ใช้แบบส่วนตัว แนะนำให้บอกรหัสผ่านหลักกับคนในครอบครัวที่ไว้วางใจสูงสุด หรือเขียนลงในซองปิดผนึกแล้วเก็บไว้ในตู้นิรภัย
  \item 1Password ยังมี "รหัสกู้คืน" ให้ด้วย เมื่อเปิดใช้งานการยืนยันตัวตนสองขั้นตอน ต้องพิมพ์บันทึกไว้เป็นอย่างยิ่ง
\end{enumerate}

\begin{tcbpractice}{ตั้งค่าตอนนี้: เตรียม "กุญแจสำรอง" ไว้ให้คนในครอบครัว}
\begin{enumerate}
  \item เลือกคนในครอบครัวที่ไว้วางใจหนึ่งคนเป็น "ผู้ติดต่อฉุกเฉิน"
  \item ตั้งค่าการเข้าถึงฉุกเฉินในตัวจัดการรหัสผ่าน (Bitwarden) หรือเข้าร่วมตู้นิรภัยที่แชร์กันในครอบครัว (1Password)
  \item เขียนรหัสผ่านหลัก + รหัสกู้คืนลงบนกระดาษ ใส่ซองปิดผนึก เก็บไว้ในตู้นิรภัยที่บ้านหรือบอกคนในครอบครัวที่เชื่อถือได้
  \item ทบทวนทุกๆ ครึ่งปี เพื่อยืนยันว่าผู้ติดต่อฉุกเฉินยังเหมาะสมอยู่หรือไม่
\end{enumerate}

\end{tcbpractice}

\booksection{แนวทางปฏิบัติที่ดีที่สุดในการใช้ตัวจัดการรหัสผ่าน}

\begin{itemize}
  \item \textbf{ลองเลย: เริ่มใช้ตัวจัดการรหัสผ่าน}เลือกตัวจัดการรหัสผ่านสักตัว (แนะนำ Bitwarden หรือแอปที่มีให้ในระบบ)
  \item \textbf{สมัครบัญชี แล้วตั้งรหัสผ่านหลักที่แข็งแรง}นำเข้ารหัสผ่านเดิมที่มีอยู่ (ถ้ามี) หรือเริ่มเพิ่มรหัสผ่านใหม่
  \item \textbf{ติดตั้งปลั๊กอินเบราว์เซอร์หรือแอปบนมือถือ แล้วลองใช้ฟีเจอร์เติมอัตโนมัติ}เช็กรายงานตรวจความปลอดภัย เพื่อหารหัสผ่านที่ควรปรับปรุง
  \item \textbf{อย่าลืมรหัสผ่านหลัก}: รหัสผ่านหลักกู้คืนไม่ได้ ลืมเมื่อไรแปลว่ารหัสผ่านทั้งหมดหายตามไปด้วย แนะนำให้จดลงกระดาษ เก็บไว้ที่ปลอดภัย
  \item \textbf{การเข้าถึงฉุกเฉิน}: ตั้งค่าผู้ติดต่อฉุกเฉินไว้ หากวันหนึ่งคุณเข้าไม่ได้ คนที่คุณไว้วางใจจะสามารถเข้าถึงรหัสผ่านของคุณได้
\end{itemize}

\begin{tcbwarn}{คำเตือนสำคัญ}
ตัวจัดการรหัสผ่านสะดวกจริง แต่ก็มีความเสี่ยง: ถ้ารหัสผ่านหลักรั่วไหล รหัสผ่านทั้งหมดก็ถูกเปิดโปง ดังนั้นรหัสผ่านหลักห้ามใช้ซ้ำกับรหัสผ่านอื่นเด็ดขาด และห้ามเขียนไว้ในที่ที่คนหาเจอง่าย อีกอย่าง แนะนำให้เลือกผู้ให้บริการที่มีสถิติด้านความปลอดภัยดี

\end{tcbwarn}

\begin{tcbpractice}{ลองเลยตอนนี้: เริ่มใช้ตัวจัดการรหัสผ่าน}
\begin{enumerate}
  \item : เลือกตัวจัดการรหัสผ่านหนึ่งตัว (แนะนำ Bitwarden หรือตัวที่มากับระบบ)
  \item : สมัครบัญชี แล้วตั้งรหัสผ่านหลักที่แข็งแรงสักรหัส
  \item : นำเข้ารหัสผ่านที่มีอยู่ (ถ้ามี) หรือเริ่มเพิ่มรหัสผ่านใหม่
  \item : ติดตั้งปลั๊กอินเบราว์เซอร์หรือแอปมือถือ ลองฟังก์ชันกรอกอัตโนมัติ
  \item : ดูรายงานตรวจสอบความปลอดภัย หารหัสผ่านที่ต้องปรับปรุง
\end{enumerate}

\end{tcbpractice}
